% !TeX root = ../main.tex
%%%%%%%%%%%%%%%%%%%%%%%%%%%%%%%%%%%%%%%%%%%%%%%%%%%%%%%%%%%%%%%%%%%%%%
%  run mode: pdfLaTeX（兼容 XeLaTeX）
%  !Mode: "Tex:Utf-8"
%  中文摘要：定义 \abstractZh 与 \keywordsZh，由 main.tex 排版输出
%%%%%%%%%%%%%%%%%%%%%%%%%%%%%%%%%%%%%%%%%%%%%%%%%%%%%%%%%%%%%%%%%%%%%%

\abstractZh{
  随着物联网与第六代移动通信（6G）技术的深度融合，海量传感节点在实时监测场景中持续产生状态更新数据。
  在频谱与能量等网络资源受限的条件下，如何高效采集并调度这些多源数据、在保证系统稳定的同时尽可能降低数据
  传输时延与信息陈旧程度，成为物联网应用亟待解决的关键问题。信息新鲜度（Age of Information，AoI）作为刻画
  信息及时性的新兴指标，能够量化接收端所持有信息相对于最新状态的新旧程度，因而受到学术界与工业界的广泛关注。

  本文面向6G物联网中的多源状态更新场景，研究了数据采集与传输调度联合优化问题，主要工作如下：第一，针对多源
  传感节点经队列汇聚至基站的网络结构，建立了状态更新与排队传输的系统模型，将问题建模为在队列稳定约束下最小化
  系统长时平均加权AoI的离散时间优化问题；第二，将深度学习与排队论相结合，利用堆叠自编码器（Stacked
  Autoencoder，SAE）对高维状态数据进行无监督特征压缩，在不显著损失信息量前提下降低单次传输的负载；第三，
  基于李雅普诺夫漂移最小化理论设计在线调度策略，在每个时隙依据当前队列状态与AoI信息决策最优的源节点进行传输，
  在不依赖未来信道与到达统计信息的前提下实现渐近最优的稳定调度。

  仿真结果表明，与经典的轮询调度、最大AoI优先及最大队列优先等基准算法相比，本文所提方法能够显著降低系统的长时
  平均加权AoI，同时保证队列稳定；结合SAE特征压缩可进一步减小传输比特开销、提升整体时效性能。本文研究可为6G
  物联网中大规模数据采集与信息及时性保障提供理论依据与工程实现参考。
}

%这里是中文关键词
\keywordsZh{信息新鲜度；堆叠自编码器；李雅普诺夫优化；多源调度；物联网；6G}
